\begin{explanationbox}
该结论揭示了椭圆上动点对长轴直径的"张角"变化规律：
\begin{itemize}
    \item \textbf{几何直观}：由于椭圆相对于长轴具有高度的对称性，点 $P$ 离长轴越远（即纵坐标 $|y|$ 越大），其对两顶点的张角越大。短轴端点是全椭圆中"最高"的点，因此视角最大。
    \item \textbf{对比圆的性质}：在圆中（$a=b$），斜率乘积为 $-1$，顶角恒为 $90^\circ$。椭圆可看作圆沿轴向压缩，因此顶角大小随位置波动。
    \item \textbf{应用价值}：
    \begin{itemize}
        \item 快速求解分式结构的最值问题。
        \item 在解析几何压轴题中，利用斜率乘积定值简化运算。
    \end{itemize}
\end{itemize}
\end{explanationbox}